\begin{abstract}
Benchmarks for AI-inpainted document forgery, such as AIForge-Doc, annotate each edit with the bounding box of the rewritten field. We ask what such labels teach a localizer. Using the authentic source images released with the training forgeries, we derive pixel-change and glyph masks and train the same network on box, pixel-change and glyph masks, changing nothing else. Under its own box metric the box-trained U-Net scores 0.92 pixel AUPRC; against the pixels that actually changed the same predictions score 0.37, and each ground truth ranks its own model first. A second decoder (DeepLabV3+) reproduces the effect. We find that a 1--3\,px residue of the green edit marker drawn by the generation pipeline survives in over 80\,\% of Gemini forgeries: removing it lowers the box model's per-image AP from 0.99 to 0.77, and pasting it alone onto an authentic image is enough to light up the field. Retraining on marker-repaired images (one seed) removes this dependence while keeping most of the box score. A tiled YOLO11 detector trained on the same boxes is the best model under the box metric (0.97 test AUPRC), depends less on the marker and transfers better to Ideogram edits (0.87 vs.\ 0.63 AP), yet reaches only 0.37 validation AUPRC under the pixel-change truth. Box scores on this benchmark therefore measure field detection, not pixel-level tamper localization; we recommend reporting both, auditing generation residue, and releasing derived masks and probes, which we provide.
\end{abstract}

\begin{IEEEkeywords}
document forgery, image forensics, tamper localization, ground truth, shortcut learning, AI inpainting, object detection
\end{IEEEkeywords}
